%=======================================================================
%  CHAPTER 3 - NUMERICAL PROBLEMS AND SOLUTIONS
%  Every pole, zero and magnitude value printed here is recomputed by
%  verify.py ch3, which roots the polynomials and evaluates H on the
%  unit circle without reference to anything written below.
%  Grouped by METHOD, not by paper: one full worked instance, then the
%  sibling papers as a data-and-answer table.
%=======================================================================
\section*{\color{acc}Ch 3 \textperiodcentered{} Analysis of LTI Systems in the Frequency
Domain\\[2pt]
{\large\textcolor{sub}{Numerical Problems \& Solutions}}}
\vspace{-4pt}{\color{acc}\rule{\textwidth}{1.1pt}}\vspace{4pt}

{\color{sub}\footnotesize \mk{35 of the 56 questions in this chapter are the same
question}: plot the poles and zeros, then sketch the magnitude response. All of them are
below. They split into two shapes, according to whether the paper gives you a difference
equation or the pole and zero locations themselves, and the second half of the work is
identical in both.}

\lead{The one procedure, used in \S1 and \S2}
\begin{enumerate}
  \item \textbf{Get $H(z)$.} From an LCCDE, Z-transform both sides and divide. Multiply
        top and bottom by $z^{N}$ so both are polynomials in $z$, not in $z^{-1}$.
  \item \textbf{Factor.} Numerator roots are zeros, denominator roots are poles. These
        are almost always quadratics, so the quadratic formula finishes it, and complex
        roots always arrive in conjugate pairs.
  \item \textbf{Plot}: unit circle, $\circ$ for each zero, $\times$ for each pole, and
        the zeros at the origin if the numerator is the shorter polynomial.
  \item \textbf{Three anchor values}, which is all a not-to-scale sketch needs:
        $|H(1)|$ at $\omega=0$, $|H(j)|$ at $\omega=\pi/2$, $|H(-1)|$ at $\omega=\pi$.
  \item \textbf{Peak at the pole angle, dip at the zero angle.} Join smoothly, then say
        in words whether it is low-pass, high-pass, band-pass or a notch.
\end{enumerate}

\sbs{%
\lead{The geometry, in one line}
\[ |H(e^{j\omega})|=|K|\,
   \frac{\prod_k|e^{j\omega}-z_k|}{\prod_k|e^{j\omega}-p_k|} \]
\begin{itemize}
  \item Distances from the zeros multiply on top, distances from the poles divide
        underneath.
  \item A pole at radius $r$ and angle $\theta$ makes $|H|$ peak at $\omega=\theta$, and
        the peak is sharper the closer $r$ is to 1.
  \item A zero \emph{on} the circle forces $|H|$ to \textbf{exactly zero} at its angle.
\end{itemize}}{%
\lead{The checks that cost nothing}
\begin{enumerate}
  \item $|H(1)|$ is the sum of the coefficients over the sum of the denominator
        coefficients; $|H(-1)|$ is the same with alternating signs. Two lines.
  \item Complex roots must be \textbf{conjugate pairs}. A lone complex pole is an
        arithmetic slip, every time.
  \item Product of the roots of $z^2+\alpha z+\beta$ is $\beta$, sum is $-\alpha$. Check
        your factorisation against the coefficients before plotting.
  \item \mk{If any pole has $|p|\ge1$ and the system is stated causal, stop}: the ROC
        excludes the unit circle and there is no frequency response to draw.
\end{enumerate}}

\hr

\T{1 \quad Difference equation given: find $H(z)$, plot, sketch $|H|$}
{\color{sub}\footnotesize \tS{24} \yr{\textbf{\texttt{80 Ch}} \m{8}, \textbf{\texttt{72 Ash}} \m{7},
\textbf{78 Bh} \m{2+4}, \textbf{81 Bh} \m{3+7}, \textbf{80 Bh} \m{3+7}, 81 Ba \m{3+7},
76 Ash \m{2+8}, \textbf{79 Bh} \m{3+7}, 79 Ba \m{3+7}, 74 Ash \m{3+7}, 72 Ka \m{6},
\textbf{72 Ch} \m{2+4}, \textbf{70 Ch} \m{6}, \textbf{\texttt{71 Bh}} \m{6},
\texttt{73 Ma} \m{7}, \textbf{69 Ch} \m{2+2+6}, \textbf{\texttt{76 Bh}} \m{2+5},
\textbf{75 Ch} \m{2+7}, \textbf{\texttt{69 Bh}} \m{3+3+2+4}, 70 Asa \m{2+6+2},
\texttt{70 Ma} \m{8}}}

\begin{asked}{2080 Chaitra \textperiodcentered{} 2072 Ashwin \textperiodcentered{} [8/7]}
\mk{Plot Magnitude Response} (not to the scale) of the system described by difference
equation. $y[n]-0.4y[n-1]+0.25y[n-2]=x[n]+0.5x[n-1]$\
\textit{Both papers print exactly this, and \textbf{neither asks for a pole-zero plot}.}
\end{asked}

{\color{sub}\footnotesize \mk{Five of the 24 ask for the magnitude response and
nothing else}: \textbf{\texttt{80 Ch}}, \textbf{\texttt{72 Ash}}, 72 Ka,
\textbf{70 Ch} and \textbf{\texttt{71 Bh}}. The pole-zero plot is worth 2 or 3 marks in
the others and \mk{nothing in these}, so the plot still helps you draw the shape but does
not earn on its own. Every other paper in this section asks for both.}

\lead{Step 1: transform and divide}
\begin{itemize}
  \item Z-transform term by term, using $y[n-k]\leftrightarrow z^{-k}Y(z)$:
  \item[] \[ Y(z)\big(1-0.4z^{-1}+0.25z^{-2}\big)=X(z)\big(1+0.5z^{-1}\big) \]
  \item[] \[ H(z)=\frac{1+0.5z^{-1}}{1-0.4z^{-1}+0.25z^{-2}}
             =\frac{z(z+0.5)}{z^{2}-0.4z+0.25} \]
  \item Multiplying by $z^{2}/z^{2}$ is what exposes the zero at the origin. \mk{Do it
        before factoring, every time.}
\end{itemize}

\lead{Step 2: factor}
\sbs{%
\begin{itemize}
  \item \textbf{Zeros}, from the numerator $z(z+0.5)$:
  \item[] \[ z=0 \quad\text{and}\quad z=-0.5 \]
  \item \textbf{Poles}, from $z^{2}-0.4z+0.25=0$:
  \item[] \[ z=\frac{0.4\pm\sqrt{0.16-1}}{2}=0.2\pm j0.4583 \]
  \item Check: product of the poles is $0.04+0.21=0.25$, the constant term. Sum is $0.4$.
        Both agree.
\end{itemize}}{%
\begin{itemize}
  \item In polar form the pole pair is
  \item[] \[ |p|=\sqrt{0.2^{2}+0.4583^{2}}=\mathbf{0.5},\qquad
             \theta=\arctan\frac{0.4583}{0.2}=66.4^{\circ}=1.159\ \text{rad} \]
  \item $|p|=0.5<1$, so the causal system is \textbf{stable} and the frequency response
        exists.
  \item \mk{$|p|=0.5$ exactly}, because the constant term of a monic quadratic is the
        product of the roots, and here that product is $|p|^2=0.25$.
\end{itemize}}

\lead{Step 3: the three anchor values}
\begin{itemize}
  \item At $\omega=0$, put $z=1$:
        $\;|H|=\left|\dfrac{1+0.5}{1-0.4+0.25}\right|=\dfrac{1.5}{0.85}=\mathbf{1.765}$
  \item At $\omega=\pi/2$, put $z=j$:
        $\;|H|=\left|\dfrac{1-0.5j}{1+0.4j-0.25}\right|
        =\dfrac{1.118}{0.850}=\mathbf{1.315}$
  \item At $\omega=\pi$, put $z=-1$:
        $\;|H|=\left|\dfrac{1-0.5}{1+0.4+0.25}\right|=\dfrac{0.5}{1.65}=\mathbf{0.303}$
\end{itemize}

\lead{Step 4: read the shape off the plot}
\begin{itemize}
  \item The pole pair sits at $66.4^{\circ}$, so $|H|$ \textbf{peaks} near there. The
        true peak is $\mathbf{1.96}$ at $\mathbf{50.9^{\circ}}$, pulled back towards
        $\omega=0$ by the zero at $z=-0.5$ sitting on the far side.
  \item The zero at $z=-0.5$ lies at angle $180^{\circ}$, so $|H|$ is \textbf{smallest at
        $\omega=\pi$}, and $0.303$ is the minimum.
  \item The zero at the origin is equidistant from every point of the circle and
        \textbf{does not shape the response at all}. It still belongs on the plot.
  \item Verdict: a \textbf{low-pass} response, rising a little to a broad hump around
        $50^{\circ}$ before falling to $0.303$ at $\omega=\pi$.
\end{itemize}

\lead{The same poles with the zero moved: 78 Bh and 81 Bh}
{\color{sub}\footnotesize Two more papers ask the identical question with
$x[n]-0.4x[n-1]$ on the right instead of $x[n]+0.5x[n-1]$. \mk{The denominator is
unchanged, so the poles are the same $0.2\pm j0.4583$}; only the zero moves, from $-0.5$
to $+0.4$, that is from angle $180^{\circ}$ to angle $0^{\circ}$.}
\begin{itemize}
  \item The zero now sits nearest $\omega=0$, so it \textbf{pulls DC down} instead of
        $\omega=\pi$: $|H(1)|=0.6/0.85=\mathbf{0.706}$,
        $|H(j)|=\mathbf{1.267}$, $|H(-1)|=1.4/1.65=\mathbf{0.849}$.
  \item Peak $\mathbf{1.34}$ at $\mathbf{73.9^{\circ}}$, and the response is now
        \textbf{high-pass} in character.
  \item \mk{One sign in the numerator flips the filter from low-pass to high-pass, with
        the poles untouched.} That is the lesson of the whole section.
\end{itemize}

\lead{Every other paper in this group}
{\color{sub}\footnotesize Same five steps. Poles from the $y$ coefficients, zeros from
the $x$ coefficients, then the three anchor values.}
\par\smallskip
{\footnotesize
\begin{tabular}{@{}L{2.0cm}L{5.6cm}L{2.4cm}L{2.3cm}L{1.35cm}L{1.75cm}@{}}
\toprule
\hrow \thead{Papers} & \thead{Question as printed} & \thead{Zeros}
& \thead{Poles ($|p|$)} & \thead{$|H|$ at $0,\frac{\pi}{2},\pi$} & \thead{Character} \\
\textbf{80 Bh} \m{3+7} & Plot the pole-zero on the z-plane and draw
  Magnitude response (not to the scale) of an LTI system described by the equation,
  $y(n)=x(n)+0.8x(n{-}1)+0.8x(n{-}2)+0.49y(n{-}2)$.
  & $-0.4\pm j0.8$ \newline $(0.894)$ & $\pm0.7$ \newline $(0.7)$
  & $5.10$ \newline $0.55$ \newline $1.96$ & low-pass, dip at $116^{\circ}$ \\
\hrow \textbf{79 Bh} \m{3+7} & Plot the pole-zero in z-plane and draw the magnitude
  response (not to the scale) of the equation of the system describe by difference
  equation: $y[n]-0.35y[n{-}1]+0.25y[n{-}2]=x[n]-0.75x[n{-}1]$.
  & $0.75$, $0$ & $0.175\pm j0.468$ \newline $(0.5)$
  & $0.28$ \newline $1.51$ \newline $1.09$ & high-pass, peak $1.53$ at $83^{\circ}$ \\
79 Ba \m{3+7} & Plot the pole-zero in z-plane and draw magnitude response (not to the
  scale) of the system described by difference equation
  $y[n]-0.3y[n{-}1]+0.2y[n{-}2]=x[n]-0.5x[n{-}1]$
  & $0.5$, $0$ & $0.15\pm j0.421$ \newline $(0.447)$
  & $0.56$ \newline $1.31$ \newline $1.00$ & high-pass, peak $1.31$ at $85^{\circ}$ \\
\hrow 74 Ash \m{3+7} & Plot the pole-zero in z-plane and Draw Magnitude Response (not to
  the scale) of the system described by difference equation
  $y[n]-0.4y[n{-}1]+0.2y[n{-}2]=x[n]+0.5x[n{-}1]+0.6x[n{-}2]+0.8x[n{-}3]$ \newline
  {\color{sub}\mk{third order in $x$}, so three zeros, not two}
  & $0.185\pm j0.942$ \newline $(0.96)$, $-0.869$ & $0.2\pm j0.4$ \newline $(0.447)$
  & $3.63$ \newline $0.56$ \newline $0.19$ & low-pass, dip at $79^{\circ}$ \\
\bottomrule
\end{tabular}}
\par\smallskip
{\footnotesize
\begin{tabular}{@{}L{2.0cm}L{5.6cm}L{2.4cm}L{2.3cm}L{1.35cm}L{1.75cm}@{}}
\toprule
\hrow \thead{Papers} & \thead{Question as printed} & \thead{Zeros} & \thead{Poles ($|p|$)} & \thead{$|H|$ at $0,\frac{\pi}{2},\pi$} & \thead{Character} \\
72 Ka \m{6} & \mk{Plot Magnitude Response} (not to the scale) of the system described by
  difference equation. $y[n]-0.5y[n{-}1]+0.3y[n{-}2]=x[n]+0.7x[n{-}1]$ \newline
  {\color{sub}\mk{no pole-zero plot asked}: the response only}
  & $-0.7$, $0$ & $0.25\pm j0.487$ \newline $(0.548)$
  & $2.13$ \newline $1.42$ \newline $0.17$ & low-pass, peak $2.44$ at $50^{\circ}$ \\
\hrow \textbf{72 Ch} \m{2+4} & Plot the pole-zero in z-plane and Draw Magnitude Response
  (not to the scale) of the system described by difference equation.
  $y[n]-0.4y[n{-}1]+0.1y[n{-}2]=x[n]+0.6x[n{-}1]$
  & $-0.6$, $0$ & $0.2\pm j0.245$ \newline $(0.316)$
  & $2.29$ \newline $1.18$ \newline $0.27$ & low-pass \\
\textbf{70 Ch} \m{6} & \mk{Plot Magnitude Response} (not to the scale) of the system
  described by difference equation.
  $y[n]-0.3y[n{-}1]+0.225y[n{-}2]=x[n]+0.5x[n{-}1]$ \newline
  {\color{sub}\mk{no pole-zero plot asked}; and note $+0.5$ where 71 Bh prints $-0.4$}
  & $-0.5$, $0$ & $0.15\pm j0.45$ \newline $(0.474)$
  & $1.62$ \newline $1.35$ \newline $0.33$ & low-pass, peak $1.80$ at $55^{\circ}$ \\
\bottomrule
\end{tabular}}
\par\smallskip
{\footnotesize
\begin{tabular}{@{}L{2.0cm}L{5.6cm}L{2.4cm}L{2.3cm}L{1.35cm}L{1.75cm}@{}}
\toprule
\hrow \thead{Papers} & \thead{Question as printed} & \thead{Zeros}
& \thead{Poles ($|p|$)} & \thead{$|H|$ at $0,\frac{\pi}{2},\pi$} & \thead{Character} \\
\textbf{\texttt{71 Bh}} \m{6} & \mk{Plot Magnitude Response} (not to the
  scale) of the system described by difference equation.
  $y[n]-0.3y[n{-}1]+0.225y[n{-}2]=x[n]-0.4x[n{-}1]$ \newline
  {\color{sub}\mk{no pole-zero plot asked}}
  & $0.4$, $0$ & $0.15\pm j0.45$ \newline $(0.474)$
  & $0.65$ \newline $1.30$ \newline $0.92$ & high-pass, peak $1.32$ at $81^{\circ}$ \\
\hrow \texttt{73 Ma} \m{7} & Plot the pole-zero in z-plane and draw magnitude response
  (not to the scale) of the system described by difference equation
  $y[n]-0.3y[n{-}1]+0.2y[n{-}2]=x[n]+0.7x[n{-}1]$
  & $-0.7$, $0$ & $0.15\pm j0.421$ \newline $(0.447)$
  & $1.89$ \newline $1.43$ \newline $0.20$ & low-pass, peak $2.00$ at $48^{\circ}$ \\
\bottomrule
\end{tabular}}
\par\smallskip
{\footnotesize
\begin{tabular}{@{}L{2.0cm}L{5.6cm}L{2.4cm}L{2.3cm}L{1.35cm}L{1.75cm}@{}}
\toprule
\hrow \thead{Papers} & \thead{Question as printed} & \thead{Zeros} & \thead{Poles ($|p|$)} & \thead{$|H|$ at $0,\frac{\pi}{2},\pi$} & \thead{Character} \\
\textbf{69 Ch} \m{2+2+6} & Why do we need Difference Equation? Draw Pole-zero in Z-Plane
  and plot magnitude response (not to the scale) of the system described by difference
  equation $y[n]-0.4y[n{-}1]+0.2y[n{-}2]=x[n]+0.1x[n{-}1]-0.06x[n{-}2]$
  & $-0.3$, $0.2$ & $0.2\pm j0.4$ \newline $(0.447)$
  & $1.30$ \newline $1.19$ \newline $0.53$ & low-pass, peak $1.51$ at $55^{\circ}$ \\
\hrow \textbf{\texttt{76 Bh}} \m{2+5} & Show the pole-zero in z-plane and plot magnitude
  response (not to the scale) of the system described by difference equation.
  $y[n]-0.6y[n{-}1]+0.35y[n{-}2]=x[n]-0.1x[n{-}1]-0.2x[n{-}2]$
  & $0.5$, $-0.4$ & $0.3\pm j0.510$ \newline $(0.592)$
  & $0.93$ \newline $1.36$ \newline $0.46$ & band-pass, peak $1.91$ at $61^{\circ}$ \\
\textbf{75 Ch} \m{2+7} & Plot the pole-zero in z-plane and draw magnitude response (not
  to scale) of the system described by \mk{differential} equation
  $y(n)-0.3y(n{-}1)=2x(n{-}2)+0.7x(n{-}1)+4x(n)$ \newline
  {\color{sub}"differential" for "difference", and \mk{the $x$ terms are printed
  backwards}, $n{-}2$ first}
  & $-0.0875\pm j0.702$ \newline $(0.707)$ & $0.3$ \newline $(0.3)$
  & $9.57$ \newline $2.03$ \newline $4.08$ & low-pass, dip at $99^{\circ}$ \\
\hrow \textbf{\texttt{69 Bh}} \m{3+3+2+4} & Explain how the poles and zeros of $H(z)$
  affect the stability and the gain response of a system. Given $H(z)$ for a system with
  the difference equation $y[n]=x[n]+1.21x[n{-}2]-0.8y[n{-}1]$, plot its poles and zeros
  on the Z-plane, \mk{determine whether it is causal and stable} and sketch its gain
  response.
  & $\pm j1.1$ \newline $(1.1)$ & $-0.8$ \newline $(0.8)$
  & $1.23$ \newline $0.16$ \newline $11.05$ & high-pass, deep dip at $90^{\circ}$ \\
70 Asa \m{2+6+2}, \texttt{70 Ma} \m{8} & Given $H(z)$ for a system with the following
  difference equation: $y(n)=x(n)+x(n{-}2)$. Plot its poles and zeros in Z plane.
  Determine its magnitude response. \mk{Also, determine whether system is causal and
  stable.} \newline {\color{sub}FIR, no feedback}
  & $\pm j$ \newline \mk{on the circle} & none \newline (all-zero)
  & $2.00$ \newline $\mathbf{0}$ \newline $2.00$ & $|H|=2|\cos\omega|$, exact null \\
\bottomrule
\end{tabular}}

\creamq{75 Ch prints $2x(n-2)+0.7x(n-1)+4x(n)$ in that order}{Read the delays, not the
position on the page. The constant term is $4$, so $H(z)=(4+0.7z^{-1}+2z^{-2})/(1-0.3z^{-1})$
and $|H(1)|=6.7/0.7=9.57$. Writing it as $2+0.7z^{-1}+4z^{-2}$ instead gives $9.57$ as
well, because the coefficients still sum to $6.7$, \mk{but every other value is wrong}
and the zeros move from $-0.0875\pm j0.702$ to $-0.0875\pm j0.351$.}

\creamq{70 Asa and 70 Ma also ask whether the system is causal and stable}{Yes to both,
and the reason is one line: $y(n)=x(n)+x(n-2)$ uses only present and past inputs, so it
is causal; its impulse response is $h=\{1,0,1\}$, a finite sum, so
$\sum|h[n]|=2<\infty$ and it is stable. \mk{Every FIR system is stable}, which is why the
question is worth only 2 marks.}

\creamq{81 Ba and 76 Ash: the pole is at $z=2.75$}{$y(n)=0.67x(n)-0.3x(n-1)+2.75y(n-1)$
rearranges to $y(n)-2.75y(n-1)=0.67x(n)-0.3x(n-1)$, so
$H(z)=(0.67z-0.3)/(z-2.75)$: one zero at $z=0.448$ and \mk{one pole at $z=2.75$, far
outside the unit circle}. If the system is causal its ROC is $|z|>2.75$, which does not
contain the unit circle, so \textbf{the frequency response does not exist} and neither
the magnitude (81 Ba) nor the phase (76 Ash) can be drawn. Only the anticausal reading,
ROC $|z|<2.75$, is stable; on that reading $|H|$ climbs gently from $0.211$ at
$\omega=0$ to $0.259$ at $\omega=\pi$. \mk{State which reading you are using.} A pole at
$0.275$ rather than $2.75$ would make the causal system stable, and is the likeliest
intended value.}

\hr

\T{2 \quad Poles and zeros given directly}
{\color{sub}\footnotesize \tS{11} \yr{82 Ba \m{3+7}, 80 Ba \m{4+6}, \textbf{82 Bh} \m{2+8},
\textbf{76 Ch} \m{2+6}, 75 Ash \m{2+8}, \textbf{\texttt{81 Ch}} \m{8},
\textbf{\texttt{73 Bh}} \m{8}, \textbf{74 Ch} \m{2+8}, \textbf{\texttt{75 Bh}} \m{2+6},
\textbf{73 Ch} \m{2+8}, \textbf{71 Ch} \m{2+8}}
\gap$\cdot$\gap \mk{no algebra at all}: the factoring is done for you}

\begin{asked}{2080 Baishakh \textperiodcentered{} 2082 Bhadra \textperiodcentered{} [4+6]}
A filter has poles at $0.45\pm j1.6$ and zeros at $0.58\pm j2.06$. Map them in the
$z$-plane and plot the magnitude response (not to scale).
\end{asked}

\lead{Step 1: radii and angles, which is all the geometry needs}
\sbs{%
\begin{itemize}
  \item \textbf{Poles} $0.45\pm j1.6$:
  \item[] \[ |p|=\sqrt{0.45^{2}+1.6^{2}}=\mathbf{1.662},\qquad
             \theta_p=\arctan\frac{1.6}{0.45}=74.3^{\circ} \]
  \item \textbf{Zeros} $0.58\pm j2.06$:
  \item[] \[ |z|=\sqrt{0.58^{2}+2.06^{2}}=\mathbf{2.140},\qquad
             \theta_z=\arctan\frac{2.06}{0.58}=74.3^{\circ} \]
\end{itemize}}{%
\begin{itemize}
  \item \mk{Poles and zeros share the same angle}, $74.3^{\circ}$, and the zeros sit
        further out. So they partly cancel: the peak the poles want to make is flattened
        by the dip the zeros want to make.
  \item Both lie \textbf{outside} the unit circle, so nothing is close to the circle and
        \mk{the response is gentle, with no sharp feature anywhere}.
  \item That is the answer the 6 marks are for. The sketch is nearly flat.
\end{itemize}}

\lead{Step 2: the three anchor values}
\begin{itemize}
  \item Build $H(z)=\dfrac{(z-0.58-j2.06)(z-0.58+j2.06)}{(z-0.45-j1.6)(z-0.45+j1.6)}
        =\dfrac{z^{2}-1.16z+4.578}{z^{2}-0.9z+2.763}$, using
        (sum, product) of each conjugate pair.
  \item $|H(1)|=\left|\dfrac{1-1.16+4.578}{1-0.9+2.763}\right|
        =\dfrac{4.418}{2.863}=\mathbf{1.544}$
  \item $|H(j)|=\mathbf{1.902}$, \quad
        $|H(-1)|=\left|\dfrac{1+1.16+4.578}{1+0.9+2.763}\right|
        =\dfrac{6.738}{4.663}=\mathbf{1.446}$
  \item Peak $\mathbf{2.04}$ at $\mathbf{74^{\circ}}$, exactly the pole angle. Range
        $1.45$ to $2.04$: \mk{a shallow bump, not a resonance}.
\end{itemize}

\lead{Step 3: the phase, when asked}
{\color{sub}\footnotesize 82 Ba asks for the phase as well. Add the angle of each vector
from a zero and subtract the angle of each vector from a pole:
$\arg H=\sum_k\arg(e^{j\omega}-z_k)-\sum_k\arg(e^{j\omega}-p_k)$. With conjugate pairs
the imaginary parts cancel at $\omega=0$ and $\omega=\pi$, so \mk{the phase is zero at
both ends} and odd in between.}

\lead{The rest of the family}
\par\smallskip
{\footnotesize
\begin{tabular}{@{}L{2.6cm}L{6.0cm}L{3.0cm}L{2.8cm}L{1.35cm}L{2.2cm}@{}}
\toprule
\hrow \thead{Papers} & \thead{Question as printed} & \thead{Poles ($|p|$)}
& \thead{Zeros ($|z|$)} & \thead{$|H|$ at $0,\frac{\pi}{2},\pi$} & \thead{Shape} \\
82 Ba \m{3+7}, 80 Ba \m{4+6}, \textbf{82 Bh} \m{2+8} & Draw the poles and zeros in the
  z-plane for the system with poles at $0.45\pm j1.6$ and zeros at $0.58\pm j2.06$. Also
  plot the magnitude response (not in scale) of the system. \newline
  {\color{sub}80 Ba prefixes "Differentiate between FIR system and IIR System"; 82 Ba
  asks for the \mk{phase response as well}}
  & $0.45\pm j1.6$ \newline $(1.662)$ & $0.58\pm j2.06$ \newline $(2.140)$
  & $1.54$ \newline $1.90$ \newline $1.45$
  & shallow bump, peak $2.04$ at $74^{\circ}$ \\
\hrow \textbf{76 Ch} \m{2+6}, 75 Ash \m{2+8}, \textbf{\texttt{81 Ch}} \m{8},
  \textbf{\texttt{73 Bh}} & Draw the pole-zero in the z-plane for a system with poles at
  $0.45\pm j1.06$ and zeroes at $0.58\pm j2.06$. Also plot the magnitude response (not
  to the scale) of the system. \newline {\color{sub}\mk{$j1.06$, not $j1.6$}: one digit
  pulls the poles in and the peak goes from $2$ to $11.5$. 75 Ash omits "not to scale"}
  & $0.45\pm j1.06$ \newline $(1.152)$ & $0.58\pm j2.06$ \newline $(2.140)$
  & $3.10$ \newline $3.93$ \newline $2.09$ & \mk{real peak} $11.50$ at $67^{\circ}$ \\
\textbf{74 Ch} \m{2+8} & The poles of a system are located at: $0.45-0.77i$ and
  $-2\pm0.3i$. Map the poles and zero in the z-plane and plot the magnitude response of
  the system. \newline {\color{sub}\mk{a lone complex pole}, not a pair, and \mk{no zeros
  are given at all}}
  & $0.45-j0.77$ \newline $(0.892)$, $-2\pm j0.3$ \newline $(2.022)$
  & none & $0.20$ \newline $0.18$ \newline $0.63$ & rises towards $\omega=\pi$ \\
\hrow \textbf{\texttt{75 Bh}}, \textbf{71 Ch} \m{2+8} & The poles of a system are located
  at: $0.45+0.77i$ and $-2\pm0.3i$ and zeroes at: $1.2\pm3i$. Map the poles and zero in
  the z-plane and plot the magnitude response of the system.
  & $0.45\mp j0.77$ \newline $(0.892)$, $-2\pm j0.3$ \newline $(2.022)$
  & $1.2\pm j3$ \newline $(3.231)$
  & $1.81$ \newline $1.76$ \newline $8.76$ & high-pass, peak $8.76$ at $\pi$ \\
\textbf{73 Ch} \m{2+8} & Define the difference equation with example. The Poles of a
  system are located at:- $0.45+0.77i$ and $2\pm0.7i$ and zeros at: $1.2\pm0.43i$. Plot
  the magnitude response of this system. \newline {\color{sub}\mk{$+2$, not $-2$}, which
  is what moves its bump}
  & $0.45+j0.77$ \newline $(0.892)$, $2\pm j0.7$ \newline $(2.119)$
  & $1.2\pm j0.43$ \newline $(1.275)$ & $0.27$ \newline $0.43$ \newline $0.37$
  & broad bump at $72^{\circ}$ \\
\bottomrule
\end{tabular}}

\creamq{Every pole in this section except $0.45\pm j0.77$ lies outside the unit
circle}{$|p|$ runs $1.152$, $1.662$, $2.022$, $2.119$. So \mk{none of these systems is
both causal and stable}. The questions only ask you to map the points and sketch $|H|$,
which is legitimate as long as the ROC is taken to be an annulus that contains the unit
circle, making the system stable but \textbf{not causal}. \mk{Say so in one line and the
mark is safe.} The $0.45\pm j0.77$ pair, at $|p|=0.892$, is the only one inside.}

\creamq{74 Ch, 75 Bh, 73 Ch and 71 Ch each print a lone complex pole}{They give
$0.45-0.77i$ or $0.45+0.77i$ with \textbf{no conjugate}, alongside a pair
$-2\pm0.3i$ or $2\pm0.7i$ that is properly conjugate. \mk{A difference equation with real
coefficients cannot have an unpaired complex pole}, so the intended pole is the pair
$0.45\pm j0.77$ and the printed sign is a dropped $\pm$. Restoring the conjugate changes
the numbers, so \textbf{state which reading you used}: with the pair restored, 74 Ch has
$|H|=0.12$, $0.21$, $0.34$ at $0$, $\pi/2$, $\pi$ and peaks at $0.80$ near
$60^{\circ}$. The values tabulated above are for the question exactly as printed.}

\creamq{Where $0.45\pm j0.77$ comes from}{The local worked example these questions are
built on specifies its poles and zeros as a \emph{radius and an angle}: a zero of radius
$0.9$ at $1.0376$ rad is $0.9\cos(1.0376)+j0.9\sin(1.0376)=0.4575+j0.7751$,
\mk{which the source rounds to $0.45+j0.77$}, and a pole of radius $0.892$ at
$2.5158$ rad is $-0.728+j0.522$. \mk{Recognising the pair saves a
conversion}, and it explains why the same two decimals keep reappearing across six
different papers.}

\hr

\T{3 \quad All-pole systems, and the steady-state response to a sinusoid}
{\color{sub}\footnotesize \tP{3} \yr{\textbf{69 Bh} \m{6}, \textbf{67 Mng} \m{10},
\textbf{68 Bh} \m{7}}}

\begin{asked}{2068 Bhadra \textperiodcentered{} [7]}
For the system $y[n]-\tfrac{10}{24}y[n-1]+\tfrac{1}{24}y[n-2]=x[n]$, find the frequency
response, then determine $y[n]$ for the input
$x[n]=\sin\!\big(\tfrac{\pi}{3}n\big)+\sin\!\big(\tfrac{\pi}{5}n\big)$.
\end{asked}

\lead{Step 1: the frequency response}
\begin{itemize}
  \item There is no $x$ delay, so the numerator is $1$: the system is \textbf{all-pole}.
  \item[] \[ H(z)=\frac{1}{1-\tfrac{10}{24}z^{-1}+\tfrac{1}{24}z^{-2}}
             =\frac{z^{2}}{z^{2}-\tfrac{5}{12}z+\tfrac{1}{24}} \]
  \item Factor: $z^{2}-\tfrac{5}{12}z+\tfrac{1}{24}=(z-\tfrac14)(z-\tfrac16)$, so the
        \textbf{poles are $\tfrac14$ and $\tfrac16$}, both real, both well inside the
        circle. Two zeros sit at the origin.
  \item[] \[ H(e^{j\omega})=\frac{1}{\big(1-\tfrac14e^{-j\omega}\big)
             \big(1-\tfrac16e^{-j\omega}\big)} \]
  \item $|H|$ falls from $1.60$ at $\omega=0$ to $0.686$ at $\omega=\pi$: a gentle
        \textbf{low-pass}, with no peak because both poles are far from the circle.
\end{itemize}

\lead{Step 2: use the eigenfunction result, do not convolve}
\begin{itemize}
  \item A sinusoid in gives the same sinusoid out, scaled by $|H|$ and shifted by
        $\angle H$ at \emph{that} frequency. Evaluate $H$ twice and stop.
  \item At $\omega=\pi/3$: $\;|H|=\mathbf{1.1955}$, $\;\angle H=\mathbf{-0.3987}$ rad
        $=-22.85^{\circ}$.
  \item At $\omega=\pi/5$: $\;|H|=\mathbf{1.4159}$, $\;\angle H=\mathbf{-0.2949}$ rad
        $=-16.90^{\circ}$.
  \item[] \[ y[n]=1.1955\sin\!\Big(\frac{\pi}{3}n-0.3987\Big)
             +1.4159\sin\!\Big(\frac{\pi}{5}n-0.2949\Big) \]
  \item \mk{The lower frequency is amplified more}, which is the low-pass character
        showing up in the answer.
\end{itemize}

\lead{69 Bh: the same method, one line shorter}
\begin{itemize}
  \item $y[n]-0.8y[n-1]+0.15y[n-2]=x[n]$ gives
        $z^{2}-0.8z+0.15=(z-0.5)(z-0.3)$, so the \textbf{poles are $0.5$ and $0.3$} with
        a double zero at the origin.
  \item $|H|$: $\mathbf{2.857}$ at $\omega=0$, $0.857$ at $\pi/2$, $\mathbf{0.513}$ at
        $\pi$. Monotonically falling, so \textbf{low-pass}, peak at DC.
\end{itemize}

\creamq{67 Mng as printed has a pole exactly on the unit circle}{The paper writes
$y[n]-\tfrac56y[n-1]-\tfrac16y[n-2]-x[n]=0$. That denominator factors as
$z^{2}-\tfrac56z-\tfrac16=(z-1)(z+\tfrac16)$, putting \mk{a pole at $z=1$}. A pole on the
circle means $|H(e^{j0})|$ is infinite, the system is only marginally stable and it is
not BIBO stable, so strictly \textbf{there is no frequency response to plot}. Change the
one sign, $-\tfrac16\to+\tfrac16$, and the poles become $\tfrac12$ and $\tfrac13$, giving
an ordinary low-pass with $|H(1)|=1/[(1-\tfrac12)(1-\tfrac13)]=3$. \mk{68 Bh asks the identical question with the
$+$ sign} and poles $\tfrac14$, $\tfrac16$, which is strong evidence that the minus is a
misprint. Answer the $+$ version and note the printed one in a line.}

\hr

\T{4 \quad Zero-input, zero-state and total response}
{\color{sub}\footnotesize \tP{2} \yr{\textbf{78 Bh} \m{4}, \textbf{66 Ma} \m{10}}}

\begin{asked}{2066 Magh \textperiodcentered{} [10]}
A system is characterized by $y(n)-0.5y(n-1)-0.25y(n-2)-x(n)=0$, the input is
$x(n)=0.2^{n}u(n)$ and the initial conditions are $y(-1)=2$, $y(-2)=4$. Find (a) the zero
input response, (b) the zero state response, (c) the total response, (d) the system
function $H(z)$, (e) the poles of $H(z)$.
\end{asked}

\lead{(d) and (e) first, because everything else uses them}
\begin{itemize}
  \item[] \[ H(z)=\frac{1}{1-0.5z^{-1}-0.25z^{-2}}
             =\frac{z^{2}}{z^{2}-0.5z-0.25} \]
  \item Roots: $z=\dfrac{0.5\pm\sqrt{0.25+1}}{2}=\dfrac{1\pm\sqrt5}{4}$, that is
        \[ p_1=\mathbf{0.8090},\qquad p_2=\mathbf{-0.3090} \]
  \item Both have $|p|<1$, so the causal system is \textbf{stable}.
\end{itemize}

\lead{(a) Zero-input response: the input is switched off, the initial conditions drive it}
\begin{itemize}
  \item The form is fixed by the poles: $\;y_{zi}[n]=C_1(0.8090)^{n}+C_2(-0.3090)^{n}$.
  \item Impose $y(-1)=2$ and $y(-2)=4$ on that form and solve the pair:
  \item[] \[ C_1=1+\frac{2}{\sqrt5}=\mathbf{1.8944},\qquad
             C_2=1-\frac{2}{\sqrt5}=\mathbf{0.1056} \]
  \item[] \[ y_{zi}[n]=1.8944\,(0.8090)^{n}+0.1056\,(-0.3090)^{n},\quad n\ge0 \]
  \item Quick check by running the recurrence $y[n]=0.5y[n-1]+0.25y[n-2]$ with no input:
        $y[0]=0.5(2)+0.25(4)=2$, $y[1]=1.5$, $y[2]=1.25$, $y[3]=1$. The closed form gives
        the same four numbers.
\end{itemize}

\lead{(b) Zero-state response: the initial conditions are zero, the input drives it}
\begin{itemize}
  \item $X(z)=\dfrac{1}{1-0.2z^{-1}}$, so $Y_{zs}(z)=H(z)X(z)$ has three poles:
        $0.8090$, $-0.3090$ and $0.2$.
  \item Partial-fraction $Y_{zs}(z)/z$ and multiply back:
  \item[] \[ y_{zs}[n]=0.9612\,(0.8090)^{n}+0.1678\,(-0.3090)^{n}
             -0.1290\,(0.2)^{n},\quad n\ge0 \]
  \item The three residues sum to $1$, which must happen because $y_{zs}[0]=1$.
        \mk{That is the free check.} The last residue is exactly $-\tfrac{4}{31}$.
\end{itemize}

\lead{(c) Total response}
\begin{itemize}
  \item[] \[ y[n]=y_{zi}[n]+y_{zs}[n]
             =2.8557\,(0.8090)^{n}+0.2734\,(-0.3090)^{n}-0.1290\,(0.2)^{n} \]
  \item First few values: $3$, $2.2$, $1.89$, $1.503$, $1.2256$. Running the original
        recurrence with $y(-1)=2$, $y(-2)=4$ and $x(n)=0.2^{n}$ gives the same sequence.
  \item \mk{Natural against forced is a different split from zero-input against
        zero-state.} The natural part is everything in $(0.8090)^n$ and $(-0.3090)^n$,
        which draws from both; the forced part is only the $(0.2)^n$ term.
\end{itemize}

\creamq{78 Bh gives no initial conditions}{$y[n]-3y[n-1]-4y[n-2]=x[n]$ has
$z^{2}-3z-4=(z-4)(z+1)$, so the zero-input response is
$y_{zi}[n]=C_1(4)^{n}+C_2(-1)^{n}$ and \mk{that form is the whole answer}, since no
$y(-1)$ or $y(-2)$ is printed to pin $C_1$ and $C_2$ down. Worth adding: $|4|>1$ and
$|-1|=1$, so the zero-input response \textbf{never decays} and the system is unstable.
Say it, because the 4 marks are for recognising the form and its behaviour.}

\hr

\T{5 \quad Linear phase}
{\color{sub}\footnotesize \tP{3} \yr{70 Asa \m{2+4}, \textbf{67 Mng} \m{6},
\textbf{69 Bh} \m{6}} \gap$\cdot$\gap the derivation for general $M$ is in \S3.6}

\begin{asked}{2070 Ashad \textperiodcentered{} [2+4]}
What is the condition satisfied by a linear phase FIR filter? Show that the filter with
$h(n)=\{-1,0,1\}$ is a linear phase filter.
\end{asked}

\sbsr{0.46}{%
\lead{The condition, for the 2 marks}
\[ h[n]=\pm h[M-1-n],\qquad 0\le n\le M-1 \]
\begin{itemize}
  \item The impulse response must be \textbf{symmetric} ($+$) or \textbf{antisymmetric}
        ($-$) about its midpoint $\alpha=\dfrac{M-1}{2}$.
  \item That $\alpha$ is then the constant group delay in samples.
\end{itemize}}{%
\lead{Check the given $h(n)$}
\begin{itemize}
  \item $M=3$, so $M-1-n$ maps $0\leftrightarrow2$ and fixes $1$.
  \item $h[0]=-1$ and $h[2]=+1$, so $h[0]=-h[2]$; and $h[1]=0=-h[1]$.
  \item So $h[n]=-h[2-n]$: the sequence is \textbf{antisymmetric}, and the condition is
        met with the minus sign.
\end{itemize}}

\lead{Now show the phase really is linear}
\begin{itemize}
  \item[] \[ H(e^{j\omega})=\sum_{n=0}^{2}h[n]e^{-j\omega n}=-1+e^{-2j\omega} \]
  \item Factor out the midpoint delay $e^{-j\omega}$, which is the standard trick:
  \item[] \[ H(e^{j\omega})=e^{-j\omega}\big(e^{-j\omega}-e^{+j\omega}\big)
             =e^{-j\omega}\big(-2j\sin\omega\big)
             =2\sin\omega\;e^{-j(\omega+\pi/2)} \]
  \item[] \[ \big|H(e^{j\omega})\big|=2|\sin\omega|,\qquad
             \arg H(e^{j\omega})=-\omega-\frac{\pi}{2} \]
  \item The phase is a \textbf{straight line} in $\omega$ with slope $-1$, so
        \[ \operatorname{grd}H=-\frac{d}{d\omega}\Big(-\omega-\frac{\pi}{2}\Big)=1
           \ \text{sample}=\alpha=\frac{M-1}{2}=1 \ \checkmark \]
  \item \mk{The constant $-\pi/2$ is why this is \emph{generalized} linear phase.} It
        comes from the $j$ in $-2j\sin\omega$, that is from the antisymmetry, and it does
        not affect the group delay, so it causes no distortion.
  \item The magnitude $2|\sin\omega|$ is zero at $\omega=0$: an antisymmetric filter
        always kills DC, which is why this one is a differencer.
\end{itemize}

\creamq{67 Mng: show that $h[n]=h[N-1-n]$, $0\le n\le N-1$, is linear phase}{This is the
general case of the same argument and it is worked in full in \S3.6: pair each term with
its mirror, factor out $e^{-j\omega\alpha}$ with $\alpha=(N-1)/2$, and what is left is a
real cosine sum. \mk{Real bracket, so the whole phase is the exponential}, giving
$\arg H=-\alpha\omega$. 69 Bh asks for the same derivation restricted to odd $M$, which
is the easier case because the midpoint sample $h[\alpha]$ stands alone instead of
pairing.}
